% DRAIN OUT - Technical Approach
% Compile on Overleaf with pdfLaTeX (default settings). No external files needed.

\documentclass[11pt,a4paper]{article}

\usepackage[margin=2.2cm]{geometry}
\usepackage[T1]{fontenc}
\usepackage{lmodern}
\usepackage{microtype}
\usepackage{amsmath,amssymb}
\usepackage{booktabs,tabularx,array,longtable}
\usepackage{enumitem}
\usepackage[table]{xcolor}
\usepackage{tikz}
\usetikzlibrary{arrows.meta,positioning,fit,backgrounds,shapes.geometric}
\usepackage{titlesec}
\usepackage{fancyhdr}
\usepackage[colorlinks=true,linkcolor=accent,urlcolor=accent,citecolor=accent]{hyperref}

% ---------- Look ----------
\definecolor{accent}{HTML}{0B5CAD}
\definecolor{built}{HTML}{1E7B34}
\definecolor{planned}{HTML}{B35C00}
\definecolor{soft}{HTML}{EEF3FA}
\definecolor{ink}{HTML}{1B2330}

\titleformat{\section}{\Large\bfseries\color{accent}}{\thesection}{0.8em}{}
\titleformat{\subsection}{\large\bfseries\color{ink}}{\thesubsection}{0.7em}{}
\titleformat{\subsubsection}{\normalsize\bfseries\color{ink}}{\thesubsubsection}{0.6em}{}

\setlength{\headheight}{14pt}
\pagestyle{fancy}
\fancyhf{}
\lhead{\small\color{accent}\textbf{DRAIN OUT} \textperiodcentered{} Technical Approach}
\rhead{\small Smart India Hackathon 2026}
\cfoot{\small\thepage}
\renewcommand{\headrulewidth}{0.4pt}

\setlist{itemsep=2pt,topsep=4pt}
\renewcommand{\arraystretch}{1.25}

% Status badges: what runs today, and what is designed but not yet in the code.
\newcommand{\built}{\textcolor{built}{\small\textbf{[BUILT]}}}
\newcommand{\planned}{\textcolor{planned}{\small\textbf{[DESIGNED]}}}

\newcommand{\keybox}[2]{%
  \begin{center}\fcolorbox{accent}{soft}{\parbox{0.94\linewidth}{\vspace{4pt}%
  \textbf{\color{accent}#1}\par\smallskip #2\vspace{4pt}}}\end{center}}

% =====================================================================
\begin{document}

\begin{titlepage}
  \centering
  \vspace*{2.5cm}
  {\color{accent}\rule{\linewidth}{1.2pt}}\par\vspace{0.6cm}
  {\Huge\bfseries DRAIN OUT\par}
  \vspace{0.3cm}
  {\LARGE Technical Approach\par}
  \vspace{0.4cm}
  {\large A multi-model, street-level flood forecasting system\\ for Chennai's storm water drainage\par}
  \vspace{0.6cm}
  {\color{accent}\rule{\linewidth}{1.2pt}}\par
  \vspace{1.5cm}
  \begin{tabular}{@{}ll@{}}
    \textbf{Event}   & Smart India Hackathon 2026 \\
    \textbf{Problem} & SIH 26085 \\
    \textbf{Team}    & \emph{Your Team Name} \\
  \end{tabular}
  \vfill
  \begin{minipage}{0.85\linewidth}\small
  \textbf{Reading this document.} Every component is tagged \built{} (running in the
  prototype today, with automated tests) or \planned{} (specified here, data in hand,
  not yet coded). Numbers quoted for built components are measured on the prototype,
  not estimated.
  \end{minipage}
\end{titlepage}

\tableofcontents
\newpage

% =====================================================================
\section{Summary}

DRAIN OUT forecasts, street by street, where Chennai will flood over the next three
hours, how deep the water will get, when it will arrive and when it will clear, and
routes traffic around it. It combines two models that answer different questions:

\begin{description}[leftmargin=1.2em,style=nextline]
  \item[Model A \textendash{} Physics Nowcast Model (``usual occurring drainage prediction'') \built]
    A coupled drain-and-surface (\emph{dual drainage}) hydraulic model. Live rainfall
    runs off the terrain into the surveyed storm water drains, which carry what their
    Manning capacity allows; the excess surcharges out of real manholes and spreads
    over the streets. It answers: \emph{given the rain coming in the next three hours,
    which streets go under, when, and how deep?}
  \item[Model B \textendash{} Historical Pattern Model \planned]
    A recurrence model built from Chennai's historical flood record: flood extents for
    5, 10, 25, 50, 100 and 200-year return periods, the city flood hazard zones, the
    753 flooding points of December 2015 and 192 measured inundation points. It answers:
    \emph{how rare is this storm, and where has a storm of this rarity flooded before?}
\end{description}

A fusion layer combines the two into one risk per street, and flags where they
disagree, which is where the most useful information is (Section~\ref{sec:fusion}).

\keybox{Why two models and not one}{A single model trained or tuned on past years
carries the bias of those years: as rainfall extremes and the city itself change, its
errors drift. Model~A has no memory at all \textendash{} it is driven by physics and
the live rain, so it does not drift with the climate, but it cannot know that a street
flooded in 2015 for reasons the survey does not capture. Model~B has memory but no
timing. Each covers the other's blind spot, and Model~B is re-fitted every monsoon on a
sliding window so its recurrence statistics follow the changing climate instead of
freezing it (Section~\ref{sec:nonstationary}).}

% =====================================================================
\section{Design goals}

\begin{enumerate}[label=\textbf{G\arabic*.}]
  \item \textbf{Street resolution.} Report depth on named streets, not on grid cells or
        wards alone: a control room acts on ``Velachery Main Road, 35~cm at 16:40''.
  \item \textbf{Lead time.} A 0\textendash3~h nowcast in 5-minute steps, the window the
        problem statement asks for.
  \item \textbf{Traceability.} Every number is surveyed, derived from surveyed values, or
        a stated assumption with its value on show. No black box.
  \item \textbf{Conservation.} Water is never created or lost by the numerics; the mass
        balance is checked on every run.
  \item \textbf{Honesty about input.} A what-if scenario is never shown as a forecast;
        assumed values are marked as assumed.
  \item \textbf{Replaceable inputs.} A radar nowcast, a LiDAR terrain model or a new ward
        survey drops in without rewriting the model.
\end{enumerate}

% =====================================================================
\section{System architecture}

\begin{figure}[htbp]
\centering
\begin{tikzpicture}[
  node distance=0.55cm and 0.6cm,
  box/.style={draw=accent,fill=soft,rounded corners=3pt,align=center,
              minimum height=0.95cm,text width=3.0cm,font=\small},
  model/.style={box,fill=accent!12,text width=3.4cm,minimum height=1.25cm,font=\small\bfseries},
  plan/.style={model,draw=planned,fill=planned!10},
  out/.style={box,fill=built!10,draw=built},
  arr/.style={-{Stealth[length=2.2mm]},thick,color=ink!70}]

  % data layer
  \node[box] (rain) {Rainfall nowcast\\\footnotesize Open-Meteo, 15 min};
  \node[box,below=of rain] (dem) {Terrain (DEM)\\\footnotesize GEDTM30, 30 m};
  \node[box,below=of dem] (drains) {Drain network\\\footnotesize GCC survey + 92 ward maps};
  \node[box,below=of drains] (osm) {Streets, waterways,\\buildings\footnotesize{} (OSM)};
  \node[box,below=of osm] (hist) {Historical record\\\footnotesize return periods, 2015};

  % models
  \node[model,right=1.3cm of dem,yshift=-0.7cm] (A) {Model A\\Physics Nowcast\\\footnotesize\mdseries coupled dual drainage};
  \node[plan,right=1.3cm of osm,yshift=0.2cm] (B) {Model B\\Historical Pattern\\\footnotesize\mdseries recurrence prior};

  % fusion
  \node[plan,right=1.2cm of A,yshift=-1.55cm,text width=2.6cm] (F) {Fusion\\\footnotesize\mdseries Bayesian log-odds};

  % outputs
  \node[out,right=1.1cm of F,yshift=2.4cm] (o1) {Street flood timeline\\\footnotesize 5 min, 3 h};
  \node[out,below=of o1] (o2) {Surcharging\\manholes};
  \node[out,below=of o2] (o3) {Flood-safe routing\\\footnotesize A* / Dijkstra};
  \node[out,below=of o3] (o4) {Ward / basin\\risk table};

  \foreach \s in {rain,dem,drains,osm} \draw[arr] (\s.east) -- (A.west);
  \foreach \s in {hist,osm,rain} \draw[arr] (\s.east) -- (B.west);
  \draw[arr] (A.east) -- (F.west);
  \draw[arr] (B.east) -- (F.west);
  \foreach \o in {o1,o2,o3,o4} \draw[arr] (F.east) -- (\o.west);
  \draw[arr,dashed] (A.north east) to[bend left=18] (o1.west);

  \begin{scope}[on background layer]
    \node[draw=ink!30,dashed,rounded corners,fit=(rain)(hist),inner sep=6pt,
          label={[font=\footnotesize\bfseries,text=ink!70]above:Data layer}] {};
    \node[draw=ink!30,dashed,rounded corners,fit=(o1)(o4),inner sep=6pt,
          label={[font=\footnotesize\bfseries,text=ink!70]above:Outputs}] {};
  \end{scope}
\end{tikzpicture}
\caption{DRAIN OUT architecture. Blue: built; orange: designed. The dashed arrow is the
path that runs today: Model~A's street forecast goes to the outputs directly.}
\label{fig:arch}
\end{figure}

\subsection{Implementation status}

\begin{longtable}{@{}p{5.2cm}p{8.6cm}l@{}}
\toprule
\textbf{Component} & \textbf{What it does} & \textbf{Status}\\
\midrule\endhead
Drain network builder & Merges GCC survey and 92 digitised ward maps into one directed graph, repairs missing links & \built\\
Hydraulic capacity & Manning capacity, condition and obstruction factors, Rational-method inflow & \built\\
Terrain conditioning & Smoothing, channel burning, pit filling, survey-level correction & \built\\
Storage zones & h-minima + watershed segmentation, fill curves & \built\\
Model A (coupled model) & 1D drains + 2D surface, HEC-22 inlets, weir exchange, surcharge & \built\\
Street timeline & Floods-at, impassable-at, peak, clears-at per street, 5-min steps & \built\\
Flood-safe routing & Depth-costed A* and Dijkstra over 80,206 roads & \built\\
What-if controls & Scenario rainfall, drain condition (silting) & \built\\
Model B (historical pattern) & Return-period street signature, event rarity, non-stationary refit & \planned\\
Ensemble for Model A & Perturbed runs to turn depth into probability & \planned\\
Fusion layer & Log-odds fusion, agreement flags & \planned\\
Validation on 2015 & Hit rate, false alarms, depth error at 192 measured points & \planned\\
\bottomrule
\end{longtable}

% =====================================================================
\section{Data foundation}

Every dataset is real, named, and cross-checked against a second independent source
before use.

\begin{longtable}{@{}p{3.0cm}p{6.2cm}p{5.4cm}@{}}
\toprule
\textbf{Layer} & \textbf{Source and size} & \textbf{Cross-checked against}\\
\midrule\endhead
Drain network (ward maps) & GCC / SECON-JBA base-map sheets, 92 wards, digitised by OCR: 12,566 drains, 1,010~km, manhole to manhole & Surveyors' flow arrows; manhole chain numbering; the CSV positions\\
Manholes & 7,951 manholes with drain top, invert, road-edge level, width, depth & Ground-level control points on the sheets (agree with the DEM datum to 0.3~m)\\
Drain network (city) & GCC storm water drain survey, 10,240 drains (7,472 used outside the ward maps, 1,457~km) & Ward maps where both exist\\
Terrain & GEDTM30 bare-earth DTM, 1 arc-second (about 30~m) & Copernicus COP30; supplied elevation point file; EGM2008 geoid\\
Buildings & OpenStreetMap, 373,298 footprints as built share per cell & DEM error at surveyed manholes\\
Water bodies & OpenStreetMap lakes, tanks, ponds, 917 ways & Kept by pit filling (Pallikaranai, Porur, Puzhal survive)\\
Streets & OpenStreetMap, 80,206 roads & OSM bridge tags (1,011 bridges excluded)\\
Channels & OSM waterways (574) + basin model macro/micro drains, canals, rivers (934 lines) & Stated length vs drawn geometry (median within 2\%)\\
Rainfall & Open-Meteo 15-minute precipitation & Open-Meteo hourly feed, hour for hour\\
Historical & Flood extents at 5/10/25/50/100/200-year return periods; hazard zones (7,453 polygons, 5 classes); 753 flooding points (2015); 192 inundation points with depth & Each other (Section~\ref{sec:modelB})\\
\bottomrule
\end{longtable}

% =====================================================================
\section{Drainage network construction \built}

\subsection{Two surveys, one network}
The GCC survey covers the whole city but records almost every section as square. The
SECON-JBA ward maps draw 92 wards manhole by manhole with measured sections and
levels, but exist only as PDF drawings. The sheets were georeferenced from their own
grid lines (exact to 0.1~m) and read by OCR into drains with the surveyed invert and
section at every manhole.

Where a ward-map drain lies along a CSV drain for at least $60\%$ of the CSV drain's
length (sampled every 15~m, within 20~m), the two are the same drain measured twice:
the ward-map drain replaces the CSV drain and borrows its condition, material and
street name.

\subsection{Graph topology}
The network is a directed graph $G=(V,E)$: vertices are junctions, edges are drains.
\begin{itemize}
  \item \textbf{Junctions.} Survey points within 15~m (the width of the road two drains
        meet in) snap to one junction; ward-map points join only where they coincide
        (0.5~m).
  \item \textbf{Flow direction} comes from the invert levels: water runs from the
        higher invert to the lower. Inverts outside $[-5, 100]$~m above MSL are data
        entry errors (e.g.\ 104,860) and give ``unknown'' rather than a wrong direction.
  \item \textbf{Slope} $S = (z_{\text{start}}-z_{\text{end}})/L$ is clamped to
        $[0.0005,\,0.05]$: below, the fall is within survey error; above, on ground as
        flat as Chennai's, it is a data error.
\end{itemize}

\subsection{Missing-link awareness}
Real survey networks are broken in ways a naive graph turns into wrong answers.
Three repairs make the topology honest:
\begin{enumerate}
  \item \textbf{Gap closing.} A drain end that nothing carries on from is joined to the
        nearest other drain within 30~m, the typical gap between two survey traverses.
  \item \textbf{Unsized drains.} 798 drains have no recorded size. Left unsized they
        would carry without limit \textendash{} an infinite pipe that empties a pond in
        one step. They take the median measured section ($0.88 \times 0.70$~m) and are
        marked \emph{assumed}.
  \item \textbf{Blind ends.} 4,045 drains end where no mapped channel takes over (more
        than 100~m from any canal, river or nullah, and not at the sea). Their water is
        put back on the ground there and the surface carries it on downhill: water with
        nowhere mapped to go cannot vanish. Treating them as free outfalls instead
        would report about 30\% fewer impassable streets at 30~mm/h \textendash{} the
        size of the error this repair prevents.
\end{enumerate}

\subsection{Hydraulic capacity}
Each drain's full-flow capacity is Manning's equation:
\begin{equation}
  Q_{\text{cap}} = \frac{1}{n}\,A\,R^{2/3}\,S^{1/2},\qquad R=\frac{A}{P}
\end{equation}
with $A$ the flow area, $P$ the wetted perimeter (a closed box also wets its soffit),
$S$ the slope and $n$ from the surveyed wall material:

\begin{center}
\begin{tabular}{@{}lc@{\hspace{1.5cm}}lc@{}}
\toprule
Material & $n$ & Material & $n$\\
\midrule
Concrete / RCC & 0.015 & Stone / masonry & 0.020\\
Brick & 0.017 & Earth / soil (kacha) & 0.030\\
\bottomrule
\end{tabular}
\end{center}

The capacity actually available is reduced by the surveyed state of the drain:
\begin{equation}
  Q_{\text{eff}} = Q_{\text{cap}} \cdot f_{\text{cond}} \cdot f_{\text{obs}} \cdot \delta,
  \qquad f_{\text{cond}} \in \{1.0\ \text{good},\ 0.85\ \text{unknown},\ 0.6\ \text{bad}\},\quad
  f_{\text{obs}} \in \{1.0,\ 0.8\}
\end{equation}
where $f_{\text{obs}}=0.8$ when the survey records a junction box, pole or tank standing
in the drain, and $\delta\in[0,1]$ is the user's \emph{drain condition} setting (1 as
surveyed, 0.5 half silted). Normal depth at a given inflow is solved from the same
equation, giving depth, velocity and Froude number per drain.

\subsection{Catchments and inflow}
Each drain collects a strip $w = 60$~m wide (carriageway plus the plots draining onto
it). Catchments accumulate down the graph, so a drain's demand is the runoff of
everything upstream. Inflow follows the Rational method:
\begin{equation}
  Q = \frac{C\, i\, A}{3.6\times 10^{6}}\quad [\text{m}^3/\text{s}],\qquad C = 0.75
\end{equation}
with $i$ in mm/h and $A$ in m$^2$ ($C$ from the IS\,/\,CPHEEO range 0.70\textendash0.95
for built-up areas). For the network-wide load, spill is \emph{routed}: a drain passes
on only $\min(\text{inflow}, Q_{\text{eff}})$, so water already on the street upstream is
not counted again at every drain below.

% =====================================================================
\section{Terrain and surface model \built}

\subsection{Conditioning the DEM}
A 30~m bare-earth model misreads a dense city in predictable ways; each is corrected
by a named step, with its setting chosen by checking that real water bodies survive
and artefacts do not.
\begin{enumerate}
  \item \textbf{Median filter} ($3\times3$ cells, 90~m) removes 3\textendash5~m
        cell-to-cell jumps where buildings bleed through ($5\times5$ starts to erase
        Porur lake).
  \item \textbf{Channel burning.} Cells under every mapped river, canal and nullah are
        lowered 2~m, so road embankments and culverts are not read as dams.
  \item \textbf{Pit filling.} Closed hollows smaller than 16 cells (1.5~ha), or deeper
        than 1~m and smaller than 64 cells (6~ha), are DEM noise and filled \textendash{}
        unless a mapped water body is there.
  \item \textbf{Survey correction.} At 4,523 manholes the road-edge level is known.
        The DEM stands a median 4~m high in built-up wards. The gap is fitted as
        \begin{equation}
          z_{\text{DEM}} - z_{\text{survey}} = a + b\cdot c_{\text{build}}
        \end{equation}
        with $c_{\text{build}}$ the building share within 75~m; the clutter term
        $b\cdot c$ (about 3~m from open to fully built) is removed everywhere, and each
        surveyed area's own residual is spread with a Gaussian kernel
        ($\sigma = 800$~m). The fit was validated on wards held out of it.
\end{enumerate}

\subsection{Storage zones}
Every regional minimum at least 2~cm deep (the $h$-minima transform) seeds a zone, and
a watershed gives every cell the minimum its rain runs to. Nested hollows stay separate
zones, so two dips either side of a low ridge fill on their own and join only once
water tops the ridge. The sea and the box edge form one extra zone where water leaves
the model. Each zone $k$ has a fill curve $V_k(h)$ integrated over its cells, with only
the open (unbuilt) share of each cell able to hold water (at least 10\%).

% =====================================================================
\section{Model A: Physics Nowcast Model \built}

Model~A is a \emph{dual drainage} model: the street surface (2D) and the drain
network (1D) run as one system, exchanging water every minute.

\subsection{Rainfall input}
The next three hours of rain come from Open-Meteo's 15-minute product, trimmed to
start at the current clock step and spread evenly over three 5-minute steps (the
honest limit of the input). The feed is a stand-in for an IMD Doppler radar nowcast,
which would replace one function and nothing else. In \emph{scenario} mode a constant
rate is held for three hours; every response states which mode it is.

\subsection{Inlets}
Rain on the strip beside a drain reaches it through gully gratings (closed drains) or
over its open edge (open drains), by the FHWA HEC-22 sag-inlet equations:
\begin{equation}
  Q_{\text{grate}} = \min\!\Big( C_w P d^{3/2},\; C_o A_o \sqrt{2 g d} \Big)\cdot(1-\kappa),
  \qquad C_w = 1.66,\ C_o = 0.67
\end{equation}
for a $450\times450$~mm grating ($P = 1.8$~m, $A_o = 0.10$~m$^2$) every 30~m, half
clogged ($\kappa=0.5$, HEC-22's design state and Chennai's everyday one). Open drains
take water over both edges as a side weir with the same coefficient. While rain is
falling the inlets see 5~cm of gutter flow (IRC:SP:50); standing water over an inlet
drains in at the depth it stands.

\subsection{The 1D network}
What the inlets catch is routed down the directed drain graph. Each drain accepts at
most $Q_{\text{eff}}$ and delays what it carries as a linear reservoir with its own
travel time $t = L / v_{\text{full}}$.

\subsection{The 2D surface}
Neighbouring zones exchange water over the cell faces they share as flow over a
broad-crested weir, drowned from downstream by the Villemonte (1947) reduction:
\begin{equation}
  q = \sqrt{g}\left(\tfrac{2}{3}H_u\right)^{3/2}
      \left[1-\left(\frac{H_d}{H_u}\right)^{1.5}\right]^{0.385}
  \quad[\text{m}^3/\text{s per m of crest}]
\end{equation}
where $H_u, H_d$ are the upstream and downstream heads over the crest. This is the
Rapid Flood Spreading Method (UK Environment Agency) with an explicit weir exchange in
place of instant spreading, so water travels along the terrain towards the sea instead
of stopping where it fell. Faces more than 3~m above the lowest crest between two
zones are pruned.

\subsection{Surcharge and backflow}
What a drain cannot take comes back up at the junction it was trying to enter: the
water surcharges out of that manhole onto the zone around it and is spread by the
surface from there. This is how the model pinpoints \emph{which manholes overflow,
when, and how much}.

\subsection{Time stepping and stability}
The surface and the network exchange water every $\Delta t = 60$~s; the surface takes
two sub-steps inside each. A sub-step may close at most half the level difference
between two zones, so the explicit exchange never overshoots equal levels and stays
stable without an implicit solver.

\subsection{Mass balance}
Every volume is accounted for. Over a run,
\begin{equation}
  V_{\text{runoff}} = V_{\text{outfall}} + V_{\text{sea}} + V_{\text{surface}}
                    + V_{\text{network}} + \varepsilon ,
\end{equation}
and the automated tests require $|\varepsilon| < 10^{-6}\,V_{\text{runoff}}$ on the
whole city.

\subsection{From water levels to streets}
Every OpenStreetMap road is densified to a point every 20~m (bridges excluded: the DEM
under a deck is the river). A point's depth at step $t$ is its zone's water level less
its ground. For every stretch and every named street (grouped by name \emph{and} ward:
Chennai has a ``5th Street'' in dozens of neighbourhoods) the model reports:
\begin{itemize}
  \item \textbf{floods at}: first step deeper than 5~cm (shallower is wet tarmac);
  \item \textbf{impassable at}: first step deeper than 30~cm;
  \item \textbf{peak} depth and when it occurs;
  \item \textbf{clears at}: when it drops back below 5~cm, or still wet at 3~h;
  \item its \textbf{ward, zone, locality}, and the \textbf{river basin} and nearest named
        channel its water belongs to (Adyar, Cooum, Kosasthalayar, \dots).
\end{itemize}

\subsection{Performance}
One three-hour storm over the whole city runs in about 6\textendash15~s on a single CPU
core. Each storm is run once and shared by every user asking about it (the 12 most
recent are kept); the built model is cached on disk and reloads in seconds.

% =====================================================================
\section{Model B: Historical Pattern Model \planned}
\label{sec:modelB}

\subsection{Why recurrence needs its own model}
Floods in Chennai recur: 2005, 2015, 2021, 2023. The intuition ``a big flood comes
every $N$ years'' is right in spirit but needs care. A $T$-year event has probability
$1/T$ in any one year; the chance of at least one in the next $n$ years is
\begin{equation}
  P(\geq 1 \text{ in } n \text{ years}) = 1 - \left(1-\frac{1}{T}\right)^{n}.
\end{equation}
For $T = n = 15$ this is $64.5\%$, not certainty \textendash{} and the probability does
not ``build up'' in the years without one. What \emph{does} change over the years is
$T$ itself: rainfall extremes intensify and the city paves over its storage, so a
storm that was a 50-year event in 1990 may be a 15-year event today. A model fitted
once to a fixed record inherits the bias of that record. Model~B is therefore designed
to be re-fitted, not frozen.

\subsection{Historical inputs (in hand)}
\begin{center}
\begin{tabular}{@{}lrl@{}}
\toprule
Dataset & Features & Content\\
\midrule
Flow extent, 5-year return period & 579 & polygons, Low / Moderate / High, 128.6 km$^2$\\
Flow extent, 10-year & 678 & 172.0 km$^2$\\
Flow extent, 25-year & 549 & 220.1 km$^2$\\
Flow extent, 50-year & 491 & 251.7 km$^2$\\
Flow extent, 100-year & 422 & 266.3 km$^2$\\
Flow extent, 200-year & 383 & 273.9 km$^2$\\
Flood hazard zones & 7,453 & Very Low to Very High\\
Flooding points, December 2015 & 753 & named locations, GCC zone and division\\
Inundation points with depth & 192 & measured depth and remark\\
\bottomrule
\end{tabular}
\end{center}
(The inundation depth field's unit is not stated in the source and is confirmed before
use.)

\subsection{Street recurrence signature}
For every 20~m street point $s$ already indexed by Model~A:
\begin{itemize}
  \item $T_{\min}(s)$ \textendash{} the smallest return period whose flood extent
        covers $s$ (a street inside the 5-year extent floods often; one only inside the
        200-year extent, rarely), with the extent's severity class;
  \item $H(s)$ \textendash{} the hazard-zone class at $s$;
  \item $D_{2015}(s)$ \textendash{} distance to the nearest 2015 flooding point, and
        the measured depth there where one exists.
\end{itemize}
These are computed once; they do not depend on the storm.

\subsection{How rare is this storm?}
The rain Model~A is running on is converted to a return period through
intensity\textendash duration\textendash frequency (IDF) curves: for each duration
$d \in \{1, 3, 6, 24\}$~h, annual maxima of rainfall depth are fitted with a
Generalised Extreme Value distribution
\begin{equation}
  F(x) = \exp\!\left\{-\left[1+\xi\,\frac{x-\mu}{\sigma}\right]^{-1/\xi}\right\},
\end{equation}
and the storm's return period is $T_{\text{event}} = \max_d \frac{1}{1-F_d(x_d)}$, the
rarest of its durations. Annual maxima come from IMD gridded rainfall over Chennai.

\subsection{Non-stationary refit: removing the drifting bias}
\label{sec:nonstationary}
To follow a changing climate and a changing city, the GEV location is allowed to move
in time,
\begin{equation}
  \mu(t) = \mu_0 + \mu_1\, t ,
\end{equation}
fitted by maximum likelihood on a sliding 30-year window and refitted after every
north-east monsoon. A significant $\mu_1 > 0$ means extremes are intensifying and the
same rain is less rare than the old record says; the model's return periods move with
it rather than staying anchored to the past. This is the specific answer to ``a single
model's bias changes with the years''.

\subsection{Model B output}
For a storm of return period $T_{\text{event}}$, the historical prior that street $s$
floods is
\begin{equation}
  p_B(s) = \operatorname{logit}^{-1}\!\big(\beta_0 + \beta_1\,\mathbf{1}[T_{\min}(s)\le T_{\text{event}}]
          + \beta_2\,H(s) + \beta_3\, e^{-D_{2015}(s)/\ell}\big),
\end{equation}
with $\beta$, $\ell$ fitted on the 2015 event.

% =====================================================================
\section{Fusion of the two models \planned}
\label{sec:fusion}

\subsection{Turning Model A into a probability}
Model~A gives a depth, not a probability. An ensemble of $M$ runs perturbs the
quantities with real uncertainty \textendash{} runoff coefficient $C$, drain condition
$\delta$, inlet spacing and the rainfall itself \textendash{} and
\begin{equation}
  p_A(s) = \frac{1}{M}\sum_{m=1}^{M} \mathbf{1}\big[d_m(s) > 5\ \text{cm}\big].
\end{equation}

\subsection{Log-odds fusion}
\begin{equation}
  \operatorname{logit} p(s) = w_A \operatorname{logit} p_A(s)
                            + w_B \operatorname{logit} p_B(s) + c,
\end{equation}
with $w_A, w_B, c$ fitted on held-out historical events. Because Model~A is driven by
the live rain, $w_A$ dominates inside the 3-hour window; $w_B$ carries the most weight
where the survey is thin (outside the 92 ward maps).

\subsection{Agreement flags}
The disagreement between the models is itself a finding:
\begin{center}
\begin{tabular}{@{}lll@{}}
\toprule
Model A (physics) & Model B (history) & Meaning\\
\midrule
wet & high & \textbf{Confirmed} \textendash{} act first\\
dry & high & \textbf{Watch} \textendash{} history says yes; check blocked drains\\
wet & low & \textbf{New hotspot} \textendash{} terrain or land-use change, encroachment\\
dry & low & Clear\\
\bottomrule
\end{tabular}
\end{center}
The \emph{new hotspot} class is how the system notices change the historical record
has not seen yet: new construction, a filled tank, a lake losing storage.

% =====================================================================
\section{Flood-safe routing \built}

The 80,206 roads become a graph whose vertices are junctions and dead ends and whose
edges are the stretches between them. Each stretch is sampled every 30~m against the
flood forecast and costed by its deepest water $d$ over the trip window (from departure
to 30 minutes after; departure anywhere in the next three hours):
\begin{equation}
  c(e) =
  \begin{cases}
    L_e & d < 5\ \text{cm}\\[2pt]
    L_e\left(1 + \dfrac{d}{10\ \text{cm}}\right) & 5 \le d < 30\ \text{cm}\\[6pt]
    \infty & d \ge 30\ \text{cm (impassable)}
  \end{cases}
\end{equation}
so 20~cm of water makes a road three times as long, and a longer dry road wins over a
short wet one. Dijkstra's algorithm and A* both find the cheapest path; A* uses the
straight-line distance to the destination as its heuristic, which is admissible (no
road is shorter than a straight line and water only adds cost), so it returns the same
optimum while settling far fewer junctions. The route is set against the plain
shortest path to show the flooded kilometres avoided and the extra distance paid.

% =====================================================================
\section{Validation}

\subsection{Built: automated checks}
22 automated tests run on every change: Manning and normal-depth formulas, HEC-22 and
weir coefficients, synthetic two-bowl terrain (zones, fill curves, spill over a ridge,
mass conservation), the whole-city mass balance, Dijkstra and A* agreement, the drain
network's catchment accumulation, and the web layer's error handling.

\subsection{Designed: skill against December 2015}
Model~A is driven with the December 2015 rainfall and compared at the 753 flooding
points and 192 depth points:
\begin{equation}
  \text{POD} = \frac{H}{H+M},\qquad
  \text{FAR} = \frac{F}{H+F},\qquad
  \text{CSI} = \frac{H}{H+M+F},
\end{equation}
with $H$ hits, $M$ misses and $F$ false alarms, plus the depth error at measured
points. Until this is done the reliable output is the \emph{order and timing} in which
streets flood rather than exact centimetres; depths over about 1~m are usually DEM
error. The same comparison fits Model~B's $\beta$ and the fusion weights, on separate
folds.

% =====================================================================
\section{Software architecture \built}

\begin{itemize}
  \item \textbf{Backend}: Python, FastAPI. Endpoints (\texttt{app/api}) only validate a
        request and call a service; all logic lives in \texttt{app/services}, one module
        per concern (drain network, hydraulics, terrain, surface, coupled model,
        forecast, routing, rainfall). Services raise typed errors that one handler turns
        into JSON responses.
  \item \textbf{Numerics}: NumPy, SciPy (KD-trees, image filters), scikit-image
        (h-minima, watershed).
  \item \textbf{Caching}: the built network and terrain are stored on disk under a key
        hashed from every input file and the code that builds it; changing any input
        rebuilds exactly that piece.
  \item \textbf{Frontend}: Leaflet map, mobile-first, works in any phone browser without
        install; HTTPS for GPS.
  \item \textbf{Interfaces}: REST/JSON for the street forecast, drains, routing and
        wards, so a GCC control room or an IMD feed can integrate directly.
\end{itemize}

% =====================================================================
\section{Implemented resources \built}
\label{sec:resources}

This section inventories everything the prototype is built from and everything it
provides: the data, the external services, the engineering standards, the software,
the code, the interfaces and the checks. Each item is in the repository and in use.

\subsection{End-to-end drainage system pipeline}

\begin{figure}[htbp]
\centering
\begin{tikzpicture}[
  node distance=0.42cm,
  stage/.style={draw=accent,fill=accent!10,rounded corners=3pt,font=\small\bfseries,
                minimum width=2.9cm,minimum height=1.05cm,align=center,text width=2.7cm},
  step/.style={draw=ink!35,fill=white,rounded corners=2pt,font=\scriptsize,align=left,
               text width=8.1cm,minimum height=1.05cm,inner sep=4pt},
  mod/.style={font=\scriptsize\ttfamily,text=built,align=left,text width=3.3cm},
  arr/.style={-{Stealth[length=2mm]},thick,color=accent}]

  \node[stage] (s1) {1. Acquire};
  \node[step,right=0.3cm of s1] (d1) {GCC drain survey CSV (10,240 drains); 92 SECON-JBA ward
     sheets (PDF); GEDTM30 and COP30 DEMs; OSM roads, waterways, buildings, water bodies;
     basin channel layer; Open-Meteo rainfall};
  \node[mod,right=0.2cm of d1] (m1) {scripts/fetch\_dem.py\\scripts/fetch\_buildings.py\\services/osm.py};

  \node[stage,below=of s1] (s2) {2. Digitise};
  \node[step,right=0.3cm of s2] (d2) {Ward sheets rendered in tiles and read by OCR; georeferenced
     from their own UTM grid (0.1 m); 7,951 manholes with top, invert, road level, size;
     12,566 drains drawn manhole to manhole};
  \node[mod,right=0.2cm of d2] (m2) {scripts/digitise\_wards.py\\data/ward\_points.json};

  \node[stage,below=of s2] (s3) {3. Merge};
  \node[step,right=0.3cm of s3] (d3) {Ward-map drain replaces a CSV drain it overlaps for $\ge$60\%
     of its length (within 20 m) and borrows its condition, material, name; duplicates within
     10 m dropped};
  \node[mod,right=0.2cm of d3] (m3) {services/drain\_network.py};

  \node[stage,below=of s3] (s4) {4. Build graph};
  \node[step,right=0.3cm of s4] (d4) {Junctions snapped at 15 m; flow direction from inverts;
     invalid inverts rejected; slopes clamped to 1:2000\,--\,1:20; catchments accumulated
     down the directed graph};
  \node[mod,right=0.2cm of d4] (m4) {services/hydraulics.py\\services/drain\_network.py};

  \node[stage,below=of s4] (s5) {5. Repair};
  \node[step,right=0.3cm of s5] (d5) {Gaps up to 30 m closed; 798 unsized drains given the
     median measured section (marked assumed); 4,045 blind ends return water to the
     surface instead of losing it};
  \node[mod,right=0.2cm of d5] (m5) {services/drain\_network.py\\services/coupled.py};

  \node[stage,below=of s5] (s6) {6. Size};
  \node[step,right=0.3cm of s6] (d6) {Manning capacity per drain, reduced by condition,
     obstruction and the silting setting; normal depth, velocity and Froude number;
     Rational-method demand};
  \node[mod,right=0.2cm of d6] (m6) {services/hydraulics.py};

  \node[stage,below=of s6] (s7) {7. Couple};
  \node[step,right=0.3cm of s7] (d7) {Drains placed on the conditioned terrain's storage zones;
     HEC-22 inlets; 1D network and 2D surface exchange every 60 s; surcharge out of
     manholes; mass balance checked};
  \node[mod,right=0.2cm of d7] (m7) {services/surface.py\\services/coupled.py};

  \node[stage,below=of s7] (s8) {8. Report};
  \node[step,right=0.3cm of s8] (d8) {Per-street timeline, surcharging manholes, ward and basin
     tables, flood-safe routes; served as REST/JSON and drawn on the dashboard};
  \node[mod,right=0.2cm of d8] (m8) {services/forecast.py\\services/navigation.py\\app/api/*};

  \foreach \a/\b in {s1/s2,s2/s3,s3/s4,s4/s5,s5/s6,s6/s7,s7/s8} \draw[arr] (\a) -- (\b);
\end{tikzpicture}
\caption{The drainage system as built, stage by stage, with the module that implements
each stage (green).}
\label{fig:pipeline}
\end{figure}

\subsection{Data resources}

{\small
\begin{longtable}{@{}p{4.3cm}p{3.1cm}r p{4.6cm}@{}}
\toprule
\textbf{Resource (file)} & \textbf{Provider} & \textbf{Size} & \textbf{Used for}\\
\midrule\endhead
\texttt{gcc\_storm\_water\_drains (1).csv} & Greater Chennai Corporation & 6.7 MB & City-wide drain survey: geometry, size, invert, material, condition\\
\texttt{Ward/*.pdf} (92 sheets) & GCC / SECON-JBA, surveyed Dec 2020 & 92 files & Manhole-level drain network and levels\\
\texttt{ward\_points.json} & Digitised from the sheets & 9.4 MB & 7,951 manholes, 12,566 drains, control points\\
\texttt{ward\_sheets.json} & Sheet title blocks & 10 KB & Ward and zone index, 196 wards\\
\texttt{GEDTM30.tif} & OpenGeoHub via OpenTopography & 1 arc-sec & Bare-earth terrain (primary)\\
\texttt{COP30.tif}, Copernicus tiles & ESA Copernicus GLO-30 & 1 arc-sec & Terrain cross-check and fallback\\
\texttt{dem\_points (1).csv} & Supplied elevation points & 36 KB & Datum check of the DEM\\
\texttt{macro\_drains.csv}, \texttt{micro\_drains.csv}, \texttt{buckingham\_canal.csv}, \texttt{krishna\_water\_canal.csv}, \texttt{rivers\_streams.csv} & Chennai Basin Drainage Maps (OpenCity) & 3.2 MB & 934 named channels by river sub-basin; channel burning; basin labels\\
\texttt{chennai\_streets.geojson} & OpenStreetMap (Overpass) & 25.6 MB & 80,206 roads: street index and routing graph\\
\texttt{chennai\_waterways.geojson} & OpenStreetMap (Overpass) & 0.3 MB & 574 rivers, canals, nullahs for channel burning\\
\texttt{building\_fraction.npz} & OpenStreetMap footprints & 0.2 MB & Built share per cell (373,298 buildings)\\
\texttt{water\_bodies.npz} & OpenStreetMap & 35 KB & 917 lakes and tanks kept as storage\\
\texttt{Historical Flood/*.csv} & Chennai flood studies & 65 MB & Model B inputs (Section~\ref{sec:modelB})\\
\bottomrule
\end{longtable}}

\subsection{External services}

{\small
\begin{longtable}{@{}p{3.4cm}p{5.0cm}p{2.3cm}p{4.0cm}@{}}
\toprule
\textbf{Service} & \textbf{Purpose} & \textbf{Access} & \textbf{Module}\\
\midrule\endhead
Open-Meteo & 15-minute rainfall nowcast; hourly and daily forecast; rain grid & Free, no key & \texttt{services/rainfall.py}, \texttt{services/weather.py}\\
OpenWeather & Current conditions; rain, temperature, wind, cloud map tiles (proxied, key kept server-side) & API key & \texttt{services/weather.py}\\
OSM Overpass API & Roads, waterways, buildings, water bodies & Free & \texttt{services/osm.py}, \texttt{scripts/fetch\_buildings.py}\\
OSM Nominatim & Place search inside the Chennai box & Free, cached & \texttt{services/navigation.py}\\
OpenTopography & GEDTM30 and COP30 terrain download & API key & \texttt{scripts/fetch\_dem.py}\\
\bottomrule
\end{longtable}}

\subsection{Engineering standards and methods implemented}

{\small
\begin{longtable}{@{}p{4.2cm}p{7.0cm}p{3.5cm}@{}}
\toprule
\textbf{Standard / method} & \textbf{Implemented as} & \textbf{Module}\\
\midrule\endhead
Manning's equation & Full-flow capacity and normal depth of every drain, open and closed sections & \texttt{hydraulics.py}\\
Rational method ($Q = CiA$) & Runoff demand, $C = 0.75$ (IS / CPHEEO built-up range) & \texttt{hydraulics.py}\\
CPHEEO Storm Water Drainage Manual & Runoff coefficient range; design-storm drains where none are surveyed (50 mm/h) & \texttt{hydraulics.py}, \texttt{coupled.py}\\
FHWA HEC-22 sag inlets & Grating weir/orifice capture, 50\% clogging, side-weir open drains & \texttt{coupled.py}\\
IRC:SP:50 & 5 cm gutter flow seen by on-grade inlets & \texttt{coupled.py}\\
Broad-crested weir + Villemonte & Zone-to-zone surface exchange, drowned flow & \texttt{coupled.py}\\
Rapid Flood Spreading Method & Storage-zone surface model (UK Environment Agency) & \texttt{surface.py}, \texttt{coupled.py}\\
Linear reservoir routing & Drain travel time $L / v_{\text{full}}$ & \texttt{coupled.py}\\
Mathematical morphology & $h$-minima, watershed, area closing for zones and pit filling & \texttt{surface.py}\\
Dijkstra, A* & Flood-costed shortest paths with admissible heuristic & \texttt{routing.py}\\
\bottomrule
\end{longtable}}

\subsection{Software stack}

{\small
\begin{longtable}{@{}p{3.2cm}p{2.0cm}p{9.5cm}@{}}
\toprule
\textbf{Component} & \textbf{Version} & \textbf{Role}\\
\midrule\endhead
Python & 3.12 & Language for the backend, models and data tools\\
FastAPI & 0.115.0 & REST API and page serving; typed request validation\\
Uvicorn & 0.30.6 & ASGI server; HTTPS for GPS on phones\\
HTTPX & 0.27.2 & Async client for every upstream API\\
NumPy & 2.5.2 & Grids, fill curves, vectorised coupled model\\
SciPy & 1.18.1 & KD-trees for snapping and search; image filters\\
scikit-image & 0.26.0 & $h$-minima, watershed, area closing\\
tifffile & 2026.8.23 & GeoTIFF DEM reading without GDAL\\
Pillow & 12.3.0 & Shaded-relief elevation image\\
Jinja2 & 3.1.6 & HTML templates\\
python-dotenv & 1.0.1 & Settings from \texttt{.env}\\
cryptography & 50.0.1 & Local HTTPS certificate\\
PyMuPDF, RapidOCR & -- & Ward-sheet rendering and OCR (data preparation only)\\
Leaflet & 1.9.4 & Interactive map in the browser\\
pytest & 9.1.1 & Automated test suite\\
\bottomrule
\end{longtable}}

\subsection{Codebase}

The backend is layered: \textbf{API} (request in, JSON out) $\rightarrow$
\textbf{services} (all logic) $\rightarrow$ \textbf{config} (paths and constants,
defined once). No endpoint contains model logic; no service depends on the web layer.

{\small
\begin{longtable}{@{}p{4.4cm}r p{9.6cm}@{}}
\toprule
\textbf{Module} & \textbf{Lines} & \textbf{Responsibility}\\
\midrule\endhead
\multicolumn{3}{@{}l}{\textit{Models and data services (\texttt{app/services})}}\\
\texttt{drain\_network.py} & 650 & Survey + ward-map merge, graph, gap repair, drain reports\\
\texttt{coupled.py} & 641 & Dual-drainage model: inlets, network, surface exchange, surcharge\\
\texttt{surface.py} & 464 & DEM conditioning, survey correction, storage zones, street depths\\
\texttt{hydraulics.py} & 452 & Manning, normal depth, Rational method, drain graph routing\\
\texttt{forecast.py} & 329 & Storm runs, cached; street, basin and manhole tables\\
\texttt{weather.py} & 313 & OpenWeather and Open-Meteo clients\\
\texttt{terrain.py} & 305 & DEM reading, grids, point lookups, relief image\\
\texttt{routing.py} & 264 & Street graph, Dijkstra, A*, depth costs\\
\texttt{navigation.py} & 165 & Geocoding, trip window, route comparison\\
\texttt{channels.py} & 136 & Basin channel layer, nearest channel and sub-basin\\
\texttt{osm.py} & 106 & Overpass downloads of roads and waterways\\
\texttt{rainfall.py} & 103 & 3-hour rainfall in 5-minute steps; scenario mode\\
\texttt{wards.py} & 89 & Ward index, zone labels, base-map files\\
\texttt{street\_index.py} & 59 & Streets densified to 20 m points\\
\texttt{diskcache.py} & 55 & Input-hashed on-disk cache of built models\\
\multicolumn{3}{@{}l}{\textit{Web layer (\texttt{app})}}\\
\texttt{api/*.py} (6 routers) & 277 & Weather, streets, drains, flood, navigation, wards\\
\texttt{main.py}, \texttt{pages.py} & 101 & App assembly, start-up warm-up, HTML pages\\
\texttt{config.py}, \texttt{errors.py} & 114 & Settings; typed errors mapped to HTTP status\\
\multicolumn{3}{@{}l}{\textit{Data preparation (\texttt{scripts})}}\\
\texttt{digitise\_wards.py} & 1,016 & Ward-sheet OCR, georeferencing, network extraction\\
\texttt{fetch\_dem.py} & 118 & OpenTopography DEM download\\
\texttt{fetch\_buildings.py} & 110 & OSM buildings and water bodies onto the DEM grid\\
\texttt{make\_cert.py} & 65 & Local HTTPS certificate for GPS on phones\\
\multicolumn{3}{@{}l}{\textit{Frontend (\texttt{app/static}, \texttt{app/templates})}}\\
10 JavaScript modules & 3,184 & Map layers, drains, flood forecast, routing, weather, UI shell\\
CSS and 3 templates & 1,601 & Mobile-first layout, light and dark themes\\
\multicolumn{3}{@{}l}{\textit{Quality assurance (\texttt{tests})}}\\
13 test files & 573 & 22 tests (Section~\ref{sec:qa})\\
\bottomrule
\end{longtable}}

\subsection{REST API catalogue}

All endpoints are \texttt{GET}, return JSON unless stated, and are documented live at
\texttt{/docs} (OpenAPI).

{\small
\begin{longtable}{@{}p{6.3cm}p{8.4cm}@{}}
\toprule
\textbf{Endpoint} & \textbf{Returns}\\
\midrule\endhead
\multicolumn{2}{@{}l}{\textit{Flood forecast}}\\
\texttt{/data-collection/rain-ponding/streets} & Street timeline, surcharging manholes, basin table, mass balance; live or scenario; filter by zone or ward\\
\texttt{/data-collection/nowcast} & Rainfall for the next 3 hours at a point\\
\texttt{/data-collection/elevation} & Elevation layer bounds, range, colour stops, DEM source\\
\texttt{/data-collection/elevation.png} & Shaded-relief image (PNG)\\
\multicolumn{2}{@{}l}{\textit{Drainage}}\\
\texttt{/data-collection/drains} & Drain network as GeoJSON with capacity, slope, catchment\\
\texttt{/data-collection/drain-hydraulics} & One drain: survey record, capacity, state at a rainfall, provenance\\
\texttt{/data-collection/network-summary} & Load across a zone or ward: bands, overloaded drains, routed spill\\
\texttt{/data-collection/drain-filters} & Zones, wards, drain types, statuses, ward-map coverage\\
\multicolumn{2}{@{}l}{\textit{Navigation}}\\
\texttt{/data-collection/route} & Flood-safe and shortest routes, trip window, Dijkstra vs A* statistics\\
\texttt{/data-collection/geocode} & Places in Chennai matching a name\\
\multicolumn{2}{@{}l}{\textit{Base layers and weather}}\\
\texttt{/data-collection/streets} & Road network (GeoJSON)\\
\texttt{/data-collection/waterways} & Rivers, canals, drains (GeoJSON)\\
\texttt{/data-collection/weather} & Current conditions by district or coordinates\\
\texttt{/data-collection/forecast} & Hourly and daily rainfall, next rain spell\\
\texttt{/data-collection/rain-grid} & Hourly rain on a $7\times7$ grid for the rain-movement map\\
\texttt{/data-collection/tiles/\{layer\}/\{z\}/\{x\}/\{y\}.png} & Weather map tile (PNG, proxied)\\
\multicolumn{2}{@{}l}{\textit{Wards and pages}}\\
\texttt{/data-collection/wards} & Every ward with zone, drain figures, base-map link\\
\texttt{/resources/ward-map/\{number\}} & Ward base-map sheet (PDF)\\
\texttt{/}, \texttt{/resources} & Dashboard and ward resources pages (HTML)\\
\bottomrule
\end{longtable}}

\subsection{Dashboard}

\begin{itemize}
  \item \textbf{Map}: drains coloured by live load, flooded street stretches by depth,
        surcharging manholes, shaded-relief elevation, weather tiles, animated rain
        movement, GPS location.
  \item \textbf{Alerts}: live or scenario banner; streets going under, impassable
        counts per 5-minute step; a time slider across the 3 hours.
  \item \textbf{Streets}: sortable table of named streets with floods-at, peak, clears-at,
        ward, locality, basin; CSV download.
  \item \textbf{Route}: from/to search, departure time, safe versus shortest route,
        flooded kilometres avoided.
  \item \textbf{What-if}: scenario rainfall rate and drain condition (silting).
  \item \textbf{Resources page}: all 196 wards with zone, drain count, length, capacity,
        drains in poor condition, and each ward's base-map PDF.
  \item Mobile-first bottom-sheet layout; light and dark themes; no installation.
\end{itemize}

\subsection{Calibration parameters}
\label{sec:params}
Every assumption is a named constant in one place, with its basis, so it can be
calibrated against observed floods without touching the model code.

{\small
\begin{longtable}{@{}p{4.2cm}r p{7.3cm}@{}}
\toprule
\textbf{Parameter} & \textbf{Value} & \textbf{Basis}\\
\midrule\endhead
Runoff coefficient $C$ & 0.75 & IS / CPHEEO, dense built-up catchment\\
Catchment strip width & 60 m & Carriageway plus plots, both sides\\
Condition factor good / unknown / bad & 1.0 / 0.85 / 0.6 & Surveyed drain status\\
Obstruction factor & 0.8 & Utility or structure in the drain section\\
Grating spacing, clogging & 30 m, 50\% & HEC-22 design state; gully positions not surveyed\\
Gutter depth & 5 cm & IRC:SP:50\\
Unsurveyed-area design storm & 50 mm/h & Rational-method drains where none are surveyed\\
Outfall-to-channel distance & 100 m & Beyond it, a drain end is blind\\
Sheet-flow speed & 0.5 m/s & Flat paved ground; lag to be calibrated\\
DEM median filter & $3\times3$ & Keeps Porur lake, removes building pits\\
Channel burn depth & 2 m & Embankments not read as dams\\
Pit filling & 16 cells; 64 cells if $>$1 m & Keeps mapped water bodies\\
Survey correction spread & $\sigma = 800$ m & One manhole's reach\\
Coupling step, surface sub-steps & 60 s, 2 & Stability of explicit exchange\\
Wet / impassable depth & 5 / 30 cm & Reporting and routing thresholds\\
Route cost & $\times(1 + d/10\,\text{cm})$ & 10 cm doubles a road's cost\\
Trip window, average speed & 30 min, 20 km/h & Cross-city drive in peak traffic\\
\bottomrule
\end{longtable}}

\subsection{Quality assurance}
\label{sec:qa}

{\small
\begin{longtable}{@{}p{4.0cm}c p{9.8cm}@{}}
\toprule
\textbf{Test file} & \textbf{Tests} & \textbf{What it proves}\\
\midrule\endhead
\texttt{test\_hydraulics.py} & 2 & Manning, normal depth, Rational method, slopes, graph routing\\
\texttt{test\_coupled.py} & 2 & Weir and grate coefficients; two-bowl terrain spills and conserves mass; city run mass balance within $10^{-6}$\\
\texttt{test\_surface.py} & 1 & Zones split at ridges; pit filling keeps real hollows; depth read-back; street timeline\\
\texttt{test\_drain\_network.py} & 3 & Flow direction from inverts; catchments never shrink downstream; routed spill never exceeds rain\\
\texttt{test\_forecast.py} & 2 & Whole-city street forecast end to end; routing agrees between A* and Dijkstra\\
\texttt{test\_routing.py} & 1 & A* settles fewer junctions for the same cost; water reroutes\\
\texttt{test\_terrain.py} & 1 & DEM covers the city, no seam at 13\textdegree N, road level below surface\\
\texttt{test\_channels.py} & 1 & Channel lengths match geometry; places find their basin\\
\texttt{test\_api.py} & 3 & Pages render; service errors return JSON with correct status; validation\\
Others (4 files) & 6 & Disk cache, rainfall scenario, street densifying, ward index\\
\midrule
\textbf{Total} & \textbf{22} & 12 fast checks on every change; 10 against the real city data\\
\bottomrule
\end{longtable}}

\subsection{Deployment and operations}
\begin{itemize}
  \item \textbf{Run}: \texttt{pip install -r requirements.txt}, then
        \texttt{python app/main.py}; HTTPS is switched on automatically when a
        certificate exists.
  \item \textbf{Configuration}: API keys and options in \texttt{.env}
        (\texttt{.env.example} documents each); nothing secret reaches the browser.
  \item \textbf{Start-up}: the drain network and flood model build in the background;
        from the disk cache this takes seconds, from scratch a few minutes.
  \item \textbf{Scaling}: each storm is computed once and shared by every user; the
        stateless API can run behind a load balancer with a shared cache directory.
\end{itemize}

% =====================================================================
\section{Limitations and roadmap}

Stated plainly, because the system is meant for real use.
\begin{longtable}{@{}p{6.4cm}p{8.0cm}@{}}
\toprule
\textbf{Limitation today} & \textbf{Upgrade path}\\
\midrule\endhead
Drain water level not solved: no pressure head or backwater from a full drain below & 1D dynamic-wave solver (SWMM-class) in place of capacity routing\\
Tide and river stage at outfalls not modelled & Couple Ennore / Adyar mouth tide gauges as boundary levels\\
30~m terrain: depth of the hollow, not the kerb & LiDAR or drone DSM (replaces the terrain files only)\\
One 60~m catchment strip and 30~m inlet spacing for all drains & GCC sub-catchment polygons and gully positions\\
Rainfall 15-minute, about 10~km grid & IMD Doppler radar nowcast at 5 minutes\\
Roads treated as two-way & Add OSM \texttt{oneway} tags (right for commuters)\\
Drain sizes outside the 92 ward maps unverified & Extend ward-map digitisation to the remaining wards\\
Unexplained 2.87~m constant between sheet levels and DEM & Settle from GCC benchmark records\\
\bottomrule
\end{longtable}

% =====================================================================
\section*{References}
\addcontentsline{toc}{section}{References}
\begin{enumerate}[label={[\arabic*]},leftmargin=2em]
  \item R.~Manning, ``On the flow of water in open channels and pipes,''
        \emph{Trans.\ Inst.\ Civil Engineers of Ireland}, 1891.
  \item FHWA, \emph{Urban Drainage Design Manual}, Hydraulic Engineering Circular
        No.~22 (HEC-22), 3rd ed., 2009.
  \item G.~R.~Villemonte, ``Submerged weir discharge studies,''
        \emph{Engineering News-Record}, 1947.
  \item Environment Agency (UK), Rapid Flood Spreading Method (RFSM) for flood risk
        assessment.
  \item Indian Roads Congress, IRC:SP:50, \emph{Guidelines on Urban Drainage}.
  \item CPHEEO, \emph{Manual on Storm Water Drainage Systems}, MoHUA, Government of
        India, 2019.
  \item S.~Coles, \emph{An Introduction to Statistical Modeling of Extreme Values},
        Springer, 2001.
  \item E.~W.~Dijkstra, ``A note on two problems in connexion with graphs,''
        \emph{Numerische Mathematik}, 1959.
  \item P.~E.~Hart, N.~J.~Nilsson, B.~Raphael, ``A formal basis for the heuristic
        determination of minimum cost paths,'' \emph{IEEE Trans.\ SSC}, 1968.
  \item P.~Soille, \emph{Morphological Image Analysis}, Springer, 2003
        ($h$-minima, watershed).
  \item Greater Chennai Corporation / SECON-JBA, storm water drain ward base maps, 2020.
  \item OpenStreetMap contributors; Open-Meteo; OpenTopography (GEDTM30, COP30).
\end{enumerate}

\end{document}
